\chapter{QBF-Kodierung bedingter Unabhängigkeit}
\label{ch:semantic-independence}

AFs können unter einer Semantik mehrere zulässige
Labellings besitzen. Bedingte Unabhängigkeit bietet eine Möglichkeit, die
Beziehungen zwischen diesen alternativen Bewertungen zu untersuchen: Sie
fragt, ob sich die Bewertungen zweier Argumentmengen bei festgehaltenem
Kontext unabhängig voneinander wählen lassen.

Die Entscheidung dieser Eigenschaft ist bereits für die hier betrachteten
Semantiken rechnerisch anspruchsvoll
\parencite{bluemel:2025:independence-aba}. QBFs eignen sich dazu, 
Anforderungen an alle Ausgangslabellings und die
Existenz eines passenden rekombinierenden Labellings gemeinsam
auszudrücken.

Auf dieser Grundlage entwickelt dieses Kapitel eine einheitliche
QBF-Kodierung bedingter Unabhängigkeit für die \emph{complete}-, die
\emph{stable}- und die \emph{preferred}-Semantik. Die Kodierung
repräsentiert Labellings boolesch und erfasst ihre Übereinstimmung und
Rekombination. Anschließend wird ihre Korrektheit gezeigt.

QBF-basierte Kodierungen von Argumentationsproblemen untersuchen
\textcite{egly:woltran:2006:reasoning-qbf}. Eine labellingbasierte
QBF-Charakterisierung verschiedener Dung-Semantiken geben
\textcite{arieli:caminada:2013:qbf-af} an. Die hier verwendete
Repräsentation knüpft an diese Grundidee an und bildet die drei möglichen
Labels eines Arguments durch jeweils eigene boolesche Variablen ab.

Abschnitt~\ref{sec:sigma-independence-af-complexity} wiederholt die
bedingte Unabhängigkeit in Argumentationssystemen nach
\textcite{rienstra:2020:independence-af}. Die
Abschnitte~\ref{sec:semantics-encoding} und
\ref{sec:uniform-independence-qbf} entwickeln anschließend die Semantik-
und Unabhängigkeitskodierung und zeigen deren Korrektheit.
\section{Bedingte Unabhängigkeit in Argumentationssystemen}
\label{sec:sigma-independence-af-complexity}

Dieser Abschnitt wiederholt die Definition der bedingten Unabhängigkeit
in Argumentationssystemen nach
\textcite{rienstra:2020:independence-af} und fasst das zugehörige
Entscheidungsproblem sowie die bekannten Komplexitätsergebnisse nach
\textcite{bluemel:2025:independence-aba} zusammen.

Intuitiv beschreibt bedingte Unabhängigkeit die Möglichkeit, 
Bewertungen verschiedener Argumentmengen miteinander zu kombinieren.
Stimmen zwei \(\sigma\)-Labellings auf der bedingenden Menge \(Z\) überein, 
so sollen die Bewertung von \(X\) aus dem ersten und die Bewertung von 
\(Y\) aus dem zweiten Labelling gemeinsam in einem weiteren 
\(\sigma\)-Labelling auftreten können, ohne die gemeinsame Bewertung von \(Z\) zu verändern.
\begin{Definition}[\cite{rienstra:2020:independence-af}]
    \label{def:sigma-independence-af}
    Sei \(F =(A,R)\) ein AF, sei \(\sigma\) eine Semantik
    und seien \(X,Y,Z\subseteq A\) paarweise disjunkt. Die Menge \(X\) ist
    \emph{\(\sigma\)-unabhängig von \(Y\) gegeben \(Z\) in \(F\)},
    geschrieben
    \(
        X\perp_\sigma Y\mid Z,
    \)
    genau dann, wenn für alle
    \(\lambda_1,\lambda_2\in\Lambda_\sigma(F)\) aus
    \(\lambda_1|_Z=\lambda_2|_Z\) die Existenz eines
    \(\lambda_3\in\Lambda_\sigma(F)\) folgt, für das
    \[
        \lambda_3|_X=\lambda_1|_X,
        \lambda_3|_Y=\lambda_2|_Y,
        \lambda_3|_Z=\lambda_1|_Z=\lambda_2|_Z
    \]
    gilt. Andernfalls sind \(X\) und \(Y\) \emph{\(\sigma\)-abhängig
    gegeben \(Z\) in \(F\)} und wir schreiben
    \(X\not\perp_\sigma Y\mid Z\).
\end{Definition}

\begin{Example}
    \label{ex:abendplanung-independence}
    Wir greifen den AF \(F_{\mathrm{Ausflug}}\) aus
    Beispiel~\ref{ex:abendplanung-complete} auf und setzen
    \[
        X=\{r_1,r_2\},\qquad
        Y=\{k_1,k_2\},\qquad
        Z=\{m,d\}.
    \]
    Die Mengen \(X\) und \(Y\) enthalten damit die Restaurant- beziehungsweise
    Kinoargumente, während \(Z\) die Stadtentscheidung festhält.

    Für die \emph{complete}-Semantik betrachten wir die Labellings
    \(\lambda_1,\lambda_2,\lambda_3\in
    \Lambda_{\mathsf{co}}(F_{\mathrm{Ausflug}})\) mit
    \[
    \begin{array}{c|c|c}
        & \mathsf{in} & \mathsf{out}\\
        \hline
        \lambda_1 & \{m,r_1,k_1\} & \{d,r_2,k_2\}\\
        \lambda_2 & \{m,r_2,k_2\} & \{d,r_1,k_1\}\\
        \lambda_3 & \{m,r_1,k_2\} & \{d,r_2,k_1\}
    \end{array}
    \]
    und jeweils
    \(\mathsf{und}(\lambda_i)=\emptyset\).
    Die Labellings \(\lambda_1\) und \(\lambda_2\) stimmen auf \(Z\)
    überein. Das Labelling \(\lambda_3\) übernimmt die Restaurantbewertung
    von \(\lambda_1\) und die Kinobewertung von \(\lambda_2\), ohne die
    Bewertung von \(Z\) zu verändern. Somit gilt
    \[
        \lambda_3|_X=\lambda_1|_X,\qquad
        \lambda_3|_Y=\lambda_2|_Y,\qquad
        \lambda_3|_Z=\lambda_1|_Z=\lambda_2|_Z.
    \]

    Diese Rekombination ist für alle auf \(Z\) übereinstimmenden
    \emph{complete}-Labellings möglich. Sind \(m\) mit
    \(\mathsf{in}\) und \(d\) mit \(\mathsf{out}\) gelabelt, können das
    Restaurant- und das Kinopaar unabhängig voneinander jeweils eine
    der Bewertungen
    \[
        (\mathsf{in},\mathsf{out}),\qquad
        (\mathsf{out},\mathsf{in}),\qquad
        (\mathsf{und},\mathsf{und})
    \]
    annehmen. Ist \(d\) mit \(\mathsf{in}\) gelabelt, sind alle Argumente
    aus \(X\cup Y\) mit \(\mathsf{out}\) gelabelt. Sind \(m\) und \(d\)
    beide mit \(\mathsf{und}\) gelabelt, gilt dies auch für alle Argumente
    aus \(X\cup Y\). Folglich gilt
    \[
        X\perp_{\mathsf{co}}Y\mid Z.
    \]

    Ohne die Festlegung der Stadt besteht diese Unabhängigkeit nicht.
    Das \emph{complete}-Labelling \(\lambda_D\) mit
    \[
        \mathsf{in}(\lambda_D)=\{d\},\qquad
        \mathsf{out}(\lambda_D)=A\setminus\{d\},\qquad
        \mathsf{und}(\lambda_D)=\emptyset
    \]
    lässt sich nicht mit der Restaurantbewertung von \(\lambda_1\)
    kombinieren: Aus \(\lambda(r_1)=\mathsf{in}\) folgt
    \(\lambda(d)=\mathsf{out}\). Dann können \(k_1\) und \(k_2\) jedoch
    nicht beide mit \(\mathsf{out}\) gelabelt sein. Daher gilt
    \[
        X\not\perp_{\mathsf{co}}Y\mid\emptyset.
    \]
\end{Example}

Bislang wurde bedingte Unabhängigkeit als semantische Eigenschaft beschrieben. 
Eine Kodierung benötigt jedoch eine konkrete Ja-Nein-Frage, 
die für eine gegebene Instanz beantwortet werden kann. 
Daher wird die Rekombinationsbedingung im nächsten Schritt als Entscheidungsproblem gefasst.
\begin{Definition}[\cite{bluemel:2025:independence-aba}]
    \label{def:sigma-independence-af-dec-problem}
    Für \(\sigma\in\{\mathsf{co},\mathsf{st},\mathsf{pr}\}\) bezeichnet
    \(\mathrm{IND}_\sigma\) das folgende \emph{Entscheidungsproblem}:
    \[
        \begin{array}{ll}
        \text{Eingabe:}
            &AF~F =(A,R)\text{ sowie paarweise disjunkte Mengen }
             X,Y,Z\subseteq A\\
        \text{Frage:}
            &\text{Gilt }X\perp_\sigma Y\mid Z\text{ in }F\,?
        \end{array}
    \]
\end{Definition}
Die Entscheidungsprobleme zur semantikspezifischen Unabhängigkeit weisen
bereits für die betrachteten Standardsemantiken eine hohe Komplexität auf.
\begin{Theorem}[\cite{bluemel:2025:independence-aba}]
\label{thm:independence-complexity}
Die Entscheidungsprobleme
\(\mathrm{IND}_{\mathsf{co}}\) und
\(\mathrm{IND}_{\mathsf{st}}\) sind
\(\Pi^{\mathrm P}_2\)-vollständig.
Das Entscheidungsproblem
\(\mathrm{IND}_{\mathsf{pr}}\) ist
\(\Pi^{\mathrm P}_3\)-vollständig.
\end{Theorem}
Die Komplexität entsteht durch die Betrachtung aller auf \(Z\) übereinstimmenden Labelling-Paare, 
für die jeweils ein geeignetes drittes Labelling existieren muss. 
Bei \emph{preferred}-Labellings erzeugt die zusätzliche Maximalitätsprüfung eine weitere Prüfebene.

Dabei ist der Randfall von AFs mit höchstens einem \(\sigma\)-Labelling zu berücksichtigen. 
Da kein Paar unterschiedlicher Bewertungen ein Gegenbeispiel liefern kann, 
gilt die Unabhängigkeit hier trivial, wie das folgende Lemma festhält.
\begin{Lemma}[\cite{rienstra:2020:independence-af}]
    \label{lem:independence-trivial-labelling-set}
    Gilt \(|\Lambda_\sigma(F)|\leq 1\), so gilt
    \(X\perp_\sigma Y\mid Z\) für alle paarweise disjunkten Mengen
    \(X,Y,Z\subseteq A\).
\end{Lemma}

\begin{proof}
    Für \(\Lambda_\sigma(F)=\emptyset\) ist die Allquantifizierung über
    \(\lambda_1\) und \(\lambda_2\) leer. Besteht
    \(\Lambda_\sigma(F)\) aus genau einem Labelling \(\lambda\), so kann
    für das einzige Paar \((\lambda,\lambda)\) dieses Labelling selbst als
    \(\lambda_3\) gewählt werden.
\end{proof}

Das Lemma ist für die spätere empirische Auswertung wesentlich: Eine hohe
Unabhängigkeitsrate kann sowohl durch viele Kombinationsmöglichkeiten als auch 
durch die Existenz von einem oder keinem Labelling entstehen.

\section{Labellingbasierte Semantikkodierung}
\label{sec:semantics-encoding}
Dieser Abschnitt entwickelt die logischen Bausteine für die spätere QBF-Kodierung bedingter Unabhängigkeit.
Dazu werden zunächst boolesche Variablen für die einzelnen Labels eingeführt und anschließend um die spezifischen Bedingungen
der \emph{stable}-, \emph{complete}- und \emph{preferred}-Semantik ergänzt.
Die daraus entstehenden Formeln bilden schließlich die Grundlage der Unabhängigkeitskodierung.

Bevor zwischen \emph{stable}-, \emph{complete}- und \emph{preferred}-Labellings unterschieden werden kann, 
muss zunächst ein beliebiges Labelling in eine boolesche Belegung übersetzt werden.
Hierzu erhält jedes Argument drei Variablen, die seine möglichen Labels repräsentieren.

\begin{Definition}
    \label{def:labelling-variable-block}
    Sei \(F=(A,R)\) ein AF. Ein
    \emph{Labellingvariablenblock} für \(F\) ist ein Tripel
    \(\mathbf{V}=(I,O,U)\) paarweise disjunkter Variablenmengen
\(
        I=\{i_a\mid a\in A\},~
        O=\{o_a\mid a\in A\},~
        U=\{u_a\mid a\in A\}.
\)
\end{Definition}
Die Variablen \(i_a\), \(o_a\) und \(u_a\) kodieren die Labels
\(\mathsf{in}\), \(\mathsf{out}\) beziehungsweise \(\mathsf{und}\)
von \(a\). 

Zusätzlich müssen wir für die Kodierung eines Labellings noch sicherstellen, dass jedes Argument 
genau ein Label erhält. 
\begin{Definition}
\label{def:One}
    Sei \(F=(A,R)\) ein AF und
    \(\mathbf V=(I,O,U)\) ein Labellingvariablenblock für \(F\).
    \[
        \operatorname{One}_F(\mathbf V)
        :=
        \bigwedge_{a\in A}
        \Bigl(
            (i_a\lor o_a\lor u_a)
            \land\neg(i_a\land o_a)
            \land\neg(i_a\land u_a)
            \land\neg(o_a\land u_a)
        \Bigr).
    \]
\end{Definition}
Die Disjunktion \(i_a\lor o_a\lor u_a\) fordert für jedes Argument
mindestens eine wahre Labelvariable. Die drei negierten Konjunktionen
schließen aus, dass mehrere Labelvariablen desselben Arguments
gleichzeitig wahr sind.

Intuitiv wählt eine Belegung, die \(\operatorname{One}_F(\mathbf V)\) erfüllt, 
für jedes Argument genau eine der drei Labelvariablen \(i_a\), \(o_a\) und \(u_a\) aus. 
Die jeweils wahre Variable legt das Label des Arguments fest. 
Somit lässt sich eine solche Belegung unmittelbar als Labelling lesen - und umgekehrt.
\begin{Definition}
    \label{def:One-model}
    Sei \(\mathbf V=(I,O,U)\) ein Labellingvariablenblock für das AF
    \(F=(A,R)\), und sei
    \(\nu\colon I\cup O\cup U\to\{0,1\}\) eine Belegung mit
    \(\nu\models\operatorname{One}_F(\mathbf V)\).
    Das durch \(\nu\) kodierte Labelling
    \(\lambda_\nu\colon A\to
    \{\mathsf{in},\mathsf{out},\mathsf{und}\}\) ist definiert durch
    \[
        \lambda_\nu(a)=
        \begin{cases}
            \mathsf{in},
                &\text{falls }\nu(i_a)=1,\\
            \mathsf{out},
                &\text{falls }\nu(o_a)=1,\\
            \mathsf{und},
                &\text{falls }\nu(u_a)=1.
        \end{cases}
    \]
\end{Definition}

\begin{Lemma}
    \label{lem:One-model}
    Sei \(F=(A,R)\) ein AF und
    \(\mathbf V=(I,O,U)\) ein Labellingvariablenblock für \(F\).
    Dann gelten die folgenden Aussagen:
    \begin{enumerate}
        \item Für jede Belegung
        \(\nu\colon I\cup O\cup U\to\{0,1\}\)
        mit
        \(\nu\models\operatorname{One}_F(\mathbf V)\)
        ist \(\lambda_\nu\) ein eindeutig bestimmtes Labelling
        von \(F\).

        \item Für jedes Labelling
        \(\lambda\colon A\to
        \{\mathsf{in},\mathsf{out},\mathsf{und}\}\)
        existiert genau eine Belegung
        \(\nu\colon I\cup O\cup U\to\{0,1\}\)
        mit
        \(\nu\models\operatorname{One}_F(\mathbf V)\)
        und
        \(\lambda_\nu=\lambda\).
    \end{enumerate}
\end{Lemma}

\begin{proof}
    Zu 1: Aus
    \(\nu\models\operatorname{One}_F(\mathbf V)\)
    folgt, dass für jedes \(a\in A\) genau eine der Variablen
    \(i_a\), \(o_a\) und \(u_a\) wahr ist. Daher trifft in
    Definition~\ref{def:One-model} genau ein Fall zu und
    \(\lambda_\nu\) ist eindeutig bestimmt.

    Zu 2: Für ein gegebenes Labelling \(\lambda\) wird eine Belegung
    \(\nu_\lambda\) definiert, indem für jedes \(a\in A\) genau die
    zu \(\lambda(a)\) gehörende Labelvariable wahr gesetzt wird.
    Dann gelten
    \(\nu_\lambda\models\operatorname{One}_F(\mathbf V)\)
    und
    \(\lambda_{\nu_\lambda}=\lambda\).

    Jede weitere Belegung mit diesen Eigenschaften muss für jedes
    Argument dieselbe Labelvariable wahr setzen und stimmt daher mit
    \(\nu_\lambda\) auf \(I\cup O\cup U\) überein. Somit ist
    \(\nu_\lambda\) eindeutig.
\end{proof}

Nach dieser Vorarbeit können wir uns der konkreten Kodierung der Semantiken stellen.
Die \emph{complete}-Semantik stellt die schwächsten Bedingungen und ist Grundlage aller weiteren hier
behandelten Semantiken. 
\begin{Definition}
    \label{def:co-st-encodings}
    Sei \(F=(A,R)\) ein AF und
    \(\mathbf V=(I,O,U)\) ein Labellingvariablenblock.
    \[
        \operatorname{Co}_F(\mathbf V)
        :=
        \operatorname{One}_F(\mathbf V)
        \land
        \bigwedge_{a\in A}
        \left(
            i_a\leftrightarrow
            \bigwedge_{b\in a^-}o_b
        \right)
        \land
        \bigwedge_{a\in A}
        \left(
            o_a\leftrightarrow
            \bigvee_{b\in a^-}i_b
        \right).
    \]
\end{Definition}
Die Teilformel \(\operatorname{One}_F(\mathbf V)\) stellt zunächst sicher, dass jedes Argument genau ein Label erhält. 
Die erste Äquivalenz weist \(a\) genau dann das Label \(\mathsf{in}\) zu, wenn alle Angreifer von \(a\) mit \(\mathsf{out}\) bezeichnet sind. 
Nach der zweiten Äquivalenz erhält \(a\) genau dann das Label \(\mathsf{out}\), wenn mindestens ein Angreifer mit \(\mathsf{in}\) bezeichnet ist. 
Damit charakterisiert \(\operatorname{Co}_F(\mathbf V)\) genau die \emph{complete}-Labellings von \(F\).

Die \emph{stable}-Semantik verschärft die Forderung der \emph{complete}-Semantik weiter.
\begin{Definition}
    Sei \(F=(A,R)\) ein AF und
    \(\mathbf V=(I,O,U)\) ein Labellingvariablenblock.
    \[
        \operatorname{St}_F(\mathbf V)
        :=
        \operatorname{Co}_F(\mathbf V)
        \land
        \bigwedge_{a\in A}\neg u_a.
    \]
\end{Definition}
Die Formel \(\operatorname{St}_F(\mathbf V)\) ergänzt die Bedingungen eines \emph{complete}-Labellings um den Ausschluss des Labels 
\(\mathsf{und}\). Sie ist daher genau für jene Belegungen erfüllt, die \emph{stable}-Labellings von \(F\) kodieren.

Im Unterschied zur \emph{complete}- und \emph{stable}-Semantik lässt sich die \emph{preferred}-Semantik nicht 
allein durch lokale Bedingungen innerhalb eines Labellings charakterisieren.
Um die geforderte Inhaltsmaximalität der In-Menge der \emph{preferred}-Semantik auszudrücken, benötigen 
wir eine zweite Kopie des Variablenblocks.
Allgemein bezeichnen wir mit \(\mathbf V^k=(I^k,O^k,U^k)\) die frische und von 
anderen Variablenblöcken disjunkte \emph{\(k\)-te Kopie} eines Variablenblocks. 
\begin{Definition}[\cite{egly:woltran:2006:reasoning-qbf}]
    \label{def:strict-inclusion}
    Seien \(\mathbf V=(I,O,U)\) und
    \(\widehat{\mathbf V}=(\widehat{I},\widehat{O},\widehat{U})\) zwei Labellingvariablenblöcke für das 
    \(AF ~F=(A,R)\).
    \[
        I < \widehat{I}
        :=
        \left(
            \bigwedge_{a\in A}(i_a\rightarrow \widehat{i}_a)
        \right)
        \land
        \left(
            \bigvee_{a\in A}(\neg i_a\land \widehat{i}_a)
        \right).
    \]
\end{Definition}
Unter einer Belegung der beiden Variablenblöcke ist \(I<\widehat I\) genau dann erfüllt, 
wenn die durch \(I\) repräsentierte In-Menge eine echte Teilmenge der durch \(\widehat I\) repräsentierten In-Menge ist. 
Die erste Konjunktion gewährleistet die Inklusion, während die zweite sicherstellt, 
dass mindestens ein weiteres Argument in \(\widehat I\) enthalten ist.

Damit haben wir alle nötigen Bausteine, um die \emph{preferred}-Semantik zu kodieren.
\begin{Definition}
    \label{def:preferred-encoding}
    Sei \(F=(A,R)\) ein AF und \(\mathbf V \text{ und } \widehat{\mathbf V}\) Labellingvariablenblöcke.
    \[
        \operatorname{Pr}_F(\mathbf V)
        :=
        \operatorname{Co}_F(\mathbf V)
        \land
        \forall\widehat{\mathbf V}\,
        \neg (\operatorname{Co}_F(\widehat{\mathbf V})
        \land I < \widehat{I}).
    \]
\end{Definition}
Die Formel \(\operatorname{Co}_F(\mathbf V)\) stellt zunächst sicher, dass \(\mathbf V\) ein \emph{complete}-Labelling kodiert. 
Der universell quantifizierte Vergleichsblock schließt anschließend jedes \emph{complete}-Labelling aus, 
dessen In-Menge die durch \(I\) bestimmte In-Menge echt enthält. Damit ist \(\operatorname{Pr}_F(\mathbf V)\) 
genau dann erfüllt, wenn das kodierte \emph{complete}-Labelling In-maximal und somit \emph{preferred} ist.

Zur Einfachheit der späteren QBF fassen wir die definierten Semantiken in 
folgender Formel zusammen.
\begin{Definition}
    \label{def:uniform-semantics-predicate}
    Für
    \(\sigma\in\{\mathsf{co},\mathsf{st},\mathsf{pr}\}\) ist
    \[
        \operatorname{Sem}^{\sigma}_F(\mathbf V)
        :=
        \begin{cases}
            \operatorname{Co}_F(\mathbf V),
                &\text{falls }\sigma=\mathsf{co},\\
            \operatorname{St}_F(\mathbf V),
                &\text{falls }\sigma=\mathsf{st},\\
            \operatorname{Pr}_F(\mathbf V),
                &\text{falls }\sigma=\mathsf{pr}.
        \end{cases}
    \]
\end{Definition}
Das Prädikat \(\operatorname{Sem}^{\sigma}_F\) fasst die drei Semantikkodierungen in einer einheitlichen Notation zusammen. 
Abhängig von \(\sigma\) bezeichnet es die Formel für \emph{complete}-, \emph{stable}- oder \emph{preferred}- Labellings. 
Dadurch können die nachfolgenden Konstruktionen unabhängig von der jeweils betrachteten Semantik formuliert werden.
\begin{Lemma}
    \label{lem:semantics-encoding-correct}
    Sei \(\nu\models\operatorname{One}_F(\mathbf V)\), und sei
    \(\lambda_\nu\) das durch \(\nu\) kodierte Labelling. Für jedes
    \(\sigma\in\{\mathsf{co},\mathsf{st},\mathsf{pr}\}\) gilt
    \[
        \nu\models\operatorname{Sem}^{\sigma}_F(\mathbf V)
        \quad\Longleftrightarrow\quad
        \lambda_\nu\in\Lambda_\sigma(F).
    \]
\end{Lemma}
\begin{proof}
    Für \(\sigma=\mathsf{co}\) geben die beiden Äquivalenzen in
    \(\operatorname{Co}_F\) genau die Definition eines \emph{complete}
    Labellings wieder. Für \(\sigma=\mathsf{st}\) kommt die Bedingung
    \(\mathsf{und}(\lambda_\nu)=\emptyset\) hinzu.

    Für \(\sigma=\mathsf{pr}\) ist
    \(\forall\widehat{\mathbf V}\,
        \neg (\operatorname{Co}_F(\widehat{\mathbf V})
        \land I < \widehat{I})\) genau dann wahr, wenn kein \emph{complete}-Labelling existiert,
    dessen In-Menge die In-Menge von \(\lambda_\nu\) echt enthält.
    Dies ist äquivalent zur In-Maximalität von \(\lambda_\nu\).
\end{proof}

\section{QBF-Kodierung der bedingten Unabhängigkeit}
\label{sec:uniform-independence-qbf}

Dieser Abschnitt entwickelt die QBF-Kodierung der bedingten Unabhängigkeit. Die zuvor eingeführten Prädikate erfassen bereits, 
ob ein einzelner Variablenblock ein gültiges Labelling der \emph{complete}-, \emph{stable}- oder \emph{preferred}-Semantik kodiert. 
Für bedingte Unabhängigkeit müssen nun drei solcher Labellings miteinander in Beziehung gesetzt werden. 
Dazu wird zunächst die Übereinstimmung zweier Labellings auf einer Teilmenge des Argumentationssystems kodiert. 
Anschließend beschreibt eine Mischungsformel die geforderte Rekombination der Bewertungen auf \(X\), \(Y\) und \(Z\). 
Beide Bausteine werden schließlich entsprechend der Quantorenstruktur der Unabhängigkeitsdefinition zu einer QBF zusammengesetzt, 
deren Korrektheit abschließend gezeigt wird.


Wir beginnen mit der Kodierung der Gleichheit zweier Labellings auf einer Teilmenge der Argumente des dazugehörigen AFs.
\begin{Definition}
    \label{def:equality-formula}
    Sei \(F =(A,R)\) ein AF, \(S\subseteq A\) und \(\mathbf V^j\) und \(\mathbf V^k\) zwei Kopien eines
    Labellingvariablenblocks.
    \[
        \operatorname{Eq}^{j,k}_S
        :=
        \bigwedge_{a\in S}
        \Bigl(
            (i_a^j\leftrightarrow i_a^k)
            \land(o_a^j\leftrightarrow o_a^k)
            \land(u_a^j\leftrightarrow u_a^k)
        \Bigr).
    \]
\end{Definition}
Jedes Argument wird in einem Variablenblock durch ein Tripel von Labelvariablen repräsentiert. Die Formel verlangt, 
dass dieses Tripel für jedes Argument aus \(S\) in beiden Blöcken identisch ist.
Kodieren die Blöcke Labellings, weisen sie den Argumenten aus \(S\) damit dieselben Labels zu.
\begin{Lemma}
    \label{lem:equality-correct}
    Sei \(\nu\) eine gemeinsame Belegung der Variablen in
    \(\mathbf V^j\) und \(\mathbf V^k\), deren Einschränkungen die
    Labellings \(\lambda_j\) und \(\lambda_k\) kodieren. Dann gilt
    \[
        \nu\models\operatorname{Eq}^{j,k}_S
        \quad\Longleftrightarrow\quad
        \lambda_j|_S=\lambda_k|_S.
    \]
\end{Lemma}
\begin{proof}
    Für jedes \(a\in S\) sind die drei Äquivalenzen genau dann erfüllt,
    wenn beide Blöcke dasselbe Label für \(a\) kodieren. Die Konjunktion
    über \(S\) liefert die Behauptung.
\end{proof}

Der zentrale Existenzschritt der bedingten Unabhängigkeit verlangt ein drittes Labelling, das auf \(X\) mit dem ersten, 
auf \(Y\) mit dem zweiten und auf \(Z\) mit dem gemeinsamen Kontext übereinstimmt. 
Wir fassen diese Anforderungen zu einem eigenständigen Baustein zusammen.
\begin{Definition}
\label{def:eq-mix-formula}
    Sei \(F =(A,R)\) ein AF, \(X,Y,Z\subseteq A\) paarweise disjunkte Mengen, 
    \(\mathbf V^1\), \(\mathbf V^2\) und \(\mathbf V^3\) paarweise disjunkte Variablenblöcke 
    und \(\operatorname{Eq}^{j,k}_S\) wie in \ref{def:equality-formula} definiert.
    \[
        \operatorname{Mix}^{1,2,3}_{X,Y\mid Z}
        :=
        \operatorname{Eq}^{1,3}_X
        \land
        \operatorname{Eq}^{2,3}_Y
        \land
        \operatorname{Eq}^{1,3}_Z.
    \]
\end{Definition}
Die Formel \( \operatorname{Mix}^{1,2,3}_{X,Y\mid Z}\) beschreibt, wie das dritte Labelling aus 
den beiden Ausgangslabellings zusammengesetzt werden soll. Es übernimmt die Bewertung von \(X\) 
aus dem ersten und die Bewertung von \(Y\) aus dem zweiten Labelling, 
während der gemeinsame Kontext \(Z\) erhalten bleibt.
\begin{Lemma}
    \label{lem:mix-correct}
    Sei \(\nu\) eine gemeinsame Belegung der Variablen in
    \(\mathbf V^1\), \(\mathbf V^2\) und \(\mathbf V^3\), deren
    Einschränkungen die Labellings
    \(\lambda_1\), \(\lambda_2\) und \(\lambda_3\) kodieren. Dann gilt
    \[
        \nu\models\operatorname{Mix}^{1,2,3}_{X,Y\mid Z}
        \quad\Longleftrightarrow\quad
        \left(
            \lambda_3|_X=\lambda_1|_X
            \land
            \lambda_3|_Y=\lambda_2|_Y
            \land
            \lambda_3|_Z=\lambda_1|_Z
        \right).
    \]
\end{Lemma}
\begin{proof}
    Die Behauptung folgt unmittelbar aus
    Lemma~\ref{lem:equality-correct} und der Definition von
    \(\operatorname{Mix}^{1,2,3}_{X,Y\mid Z}\).
\end{proof}

Damit stehen sowohl die semantische Gültigkeit der drei Labellings als
auch ihre erforderlichen Übereinstimmungen als getrennte Bausteine zur
Verfügung. Wir setzen diese nun entsprechend der Quantorenstruktur der
Unabhängigkeitsdefinition zusammen.
\begin{Definition}
    \label{def:generic-independence-qbf}
    Sei \(F=(A,R)\) ein \(AF\), \(X,Y,Z\subseteq A\) paarweise disjunkt, \(\mathbf V^1, 
    \mathbf V^2, \mathbf V^3\) paarweise disunkte Labellingvariablenblöcke und
    \(\sigma\in\{\mathsf{co},\mathsf{st},\mathsf{pr}\}\).
    \[
    \begin{split}
        \Phi^\sigma_F(X,Y\mid Z)
        :=\forall\mathbf V^1\,\forall\mathbf V^2\,
        \Bigl(&
            \operatorname{Sem}^{\sigma}_F(\mathbf V^1)
            \land
            \operatorname{Sem}^{\sigma}_F(\mathbf V^2)
            \land
            \operatorname{Eq}^{1,2}_Z
        \\
        &\rightarrow
        \exists\mathbf V^3\,
        \bigl(
            \operatorname{Sem}^{\sigma}_F(\mathbf V^3)
            \land
            \operatorname{Mix}^{1,2,3}_{X,Y\mid Z}
        \bigr)
        \Bigr).
    \end{split}
    \]
\end{Definition}
Bei der \emph{preferred}-Kodierung werden zusätzlich jeweils
paarweise disjunkte Labellingvariablenblöcke
\(\widehat{\mathbf V}^{1},
\widehat{\mathbf V}^{2},
\widehat{\mathbf V}^{3}\)
verwendet.

Die beiden universellen Variablenblöcke repräsentieren die beliebig
gewählten Labellings \(\lambda_1\) und \(\lambda_2\). Der existentielle
Block repräsentiert das kombinierende Labelling \(\lambda_3\). Die
Prämisse der Implikation beschränkt die universellen Belegungen auf
\(\sigma\)-Labellings, die auf \(Z\) übereinstimmen.
\begin{Theorem}
    \label{thm:generic-independence-qbf-correct}
    Für jedes \(AF~F=(A,R)\), alle paarweise disjunkten Mengen
    \(X,Y,Z\subseteq A\) und jedes
    \(\sigma\in\{\mathsf{co},\mathsf{st},\mathsf{pr}\}\) gilt
    \[
        \Phi^\sigma_F(X,Y\mid Z)\text{ ist wahr}
        \quad\Longleftrightarrow\quad
        X\perp_\sigma Y\mid Z \text{ in } F.
    \]
\end{Theorem}
\begin{proof}
    Sei die QBF wahr, und seien
    \(\lambda_1,\lambda_2\in\Lambda_\sigma(F)\) mit
    \(\lambda_1|_Z=\lambda_2|_Z\) gegeben. Ihre kodierenden Belegungen
    erfüllen nach Lemma~\ref{lem:semantics-encoding-correct} die beiden
    Semantikkodierungen und nach
    Lemma~\ref{lem:equality-correct} die Formel
    \(\operatorname{Eq}^{1,2}_Z\). Die QBF liefert daher eine Belegung von
    \(\mathbf V^3\), die ein \(\sigma\)-Labelling \(\lambda_3\) kodiert
    und nach Lemma~\ref{lem:mix-correct} die geforderten Restriktionen erfüllt. 
    Somit gilt \(X\perp_\sigma Y\mid Z\).

    Gelte umgekehrt \(X\perp_\sigma Y\mid Z\), und seien die beiden
    universellen Blöcke beliebig belegt. Ist die Prämisse falsch, ist die
    Implikation wahr. Andernfalls kodieren die Blöcke zwei
    \(\sigma\)-Labellings, die auf \(Z\) übereinstimmen. Nach Definition
    der Unabhängigkeit existiert ein passendes \(\lambda_3\). Seine
    kodierende Belegung erfüllt die Konklusion. Dies gilt für jede Belegung
    der universellen Blöcke, also ist die QBF wahr.
\end{proof}

Damit ist die bedingte Unabhängigkeit durch die Wahrheit von
\(\Phi_F^\sigma(X,Y\mid Z)\) charakterisiert. Die entwickelte
QBF-Kodierung bildet die Grundlage für die algorithmische Umsetzung und
die empirische Untersuchung im folgenden Kapitel.